\section{\uppercase{Conclusion}} 
\label{sec:conclusion}

This work introduced \textbf{RALA-ChangeFormer}, an efficient variant of ChangeFormer in which the multi-head self-attention modules of the \textbf{encoder} are replaced with rank-augmented linear attention. All remaining components of the architecture, including the convolutional backbone, feed-forward layers, and the lightweight MLP decoder, remain unchanged. The goal of this modification is to alleviate the quadratic cost of MHSA in high-token regimes while preserving the representational capacity required for multi-temporal building change detection.

Experimental evaluations on LEVIR-CD and DSIFN-CD show that RALA-ChangeFormer maintains segmentation accuracy comparable to the baseline model, confirming that the rank-augmented KV buffer is sufficiently expressive for bi-temporal spatial-temporal reasoning. Importantly, unlike previous linear-attention formulations, RALA avoids the low-rank collapse typically observed in dense prediction tasks, as further illustrated by the qualitative attention-map comparisons.

From a computational perspective, the proposed encoder-level integration yields consistent reductions in attention-related FLOPs and peak GPU memory usage, and increases feasible batch sizes, particularly at higher resolutions where attention costs dominate. These improvements enhance practical scalability under fixed memory budgets, enabling the model to process more patches per iteration and to support efficient training in high-resolution or large-batch settings without increasing hardware requirements. Overall, the results demonstrate that RALA-ChangeFormer provides a favorable accuracy-efficiency trade-off for VHR change detection, preserving predictive behaviour while reducing resource demands.

Rather than proposing a novel change detection formulation, this work aims to empirically assess whether modern linear-attention mechanisms can be integrated into hierarchical bi-temporal change detection architectures without degrading performance, and to quantify the resulting efficiency-accuracy trade-offs. In this context, several promising directions remain for future work. One avenue is to explore adaptive or data-dependent strategies that refine the rank selection in RALA, potentially drawing on matrix reconstruction techniques such as SVD to further enhance representational diversity. Another promising direction is the integration of RALA into multi-scale or multi-task BCD frameworks, enabling a more robust evaluation of its generalization capacity across heterogeneous settings. Finally, with generative approaches—such as diffusion and flow-matching models—emerging as powerful alternatives for BCD, a key opportunity lies in developing memory-efficient attention formulations that can be embedded into these computationally demanding architectures.